% Original PDF pp. 16; Table E1.
\begingroup
\renewcommand{\thetable}{E1}
\begin{longtable}{@{}>{\raggedright\arraybackslash}p{.24\linewidth} >{\raggedright\arraybackslash}p{.32\linewidth} >{\raggedright\arraybackslash}p{.38\linewidth}@{}}
\caption{Result placeholders. ``Not run'' is not zero failures and is not a claim that the repository has no tests.}\label{tab:E1}\\
\toprule
\textbf{Measurement} & \textbf{Planned record} & \textbf{Present status} \\
\midrule
\endfirsthead
\textbf{Measurement} & \textbf{Planned record} & \textbf{Present status} \\
\midrule
\endhead
\bottomrule
\endlastfoot
Legal source-to-IR fidelity & Expert-reviewed atoms, exceptions, applicability and source spans & Not run \\[3pt]
CVE restriction specificity & Vulnerable positives, fixed negatives, unrelated-code controls & Not run \\[3pt]
Skill grounding & Source spans, control-flow fidelity, malicious Markdown negatives & Not run \\[3pt]
Proof job execution & Per-family provider/checker route, assumptions, result and receipt & Not run \\[3pt]
Guarded effect prevention & Delegate/effect counters under allow, deny, unknown, stale and replay & Existing assertions inspected; frozen run pending \\[3pt]
DuckDB owner safety & Wrong generation/fence, concurrent commands, lost reply, restart & Existing implementation and test design inspected; run pending \\[3pt]
Cross-transport parity & Stdio/HTTP/P2P same request and outcome; actual network mode & Existing suite inspected; run pending \\[3pt]
End-to-end efficiency & All model tokens, context retrieval, proof and state/storage overhead & Not run \\[3pt]
\end{longtable}
\endgroup
